% Hafta 4 — Yığın taşması ve kanarya
% Derleme: py -3.12 tools/latex/derle.py docs/week-4/assets/h04-04-yigin-kanarya.tex
\documentclass[tikz,border=8pt]{standalone}
\usepackage{cen429-cizim}
\begin{document}
\begin{tikzpicture}[node distance=3mm and 3mm]
  % Başlık
  \node[baslik] (b) at (0,0) {Yığın taşması ve kanarya};
  \node[altbaslik] at (b.south west) {taşma tampondan dönüş adresine doğru yazar};

  % Yığın: en üstte tampon, en altta dönüş adresi
  \node[kutu=46mm, below=16mm of b.south west, anchor=north west] (buf) {char buf[8]\\(yerel tampon)};
  \node[uyarikutu=46mm, below=of buf] (kan) {\kb{KANARYA}\\(rastgele nöbetçi)};
  \node[morkutu=46mm, below=of kan] (rbp) {saved rbp};
  \node[kotukutu=46mm, below=of rbp] (don) {\kb{DÖNÜŞ ADRESİ}};

  % Taşan yazma oku: yığının solunda, aşağıdan yukarıya kalın kırmızı ok
  \coordinate (okalt) at ($(don.south west)+(-9mm,2mm)$);
  \coordinate (oküst) at ($(buf.north west)+(-9mm,-2mm)$);
  \draw[kotuok, line width=2.5pt] (okalt) -- (oküst);
  \node[text=kotu, font=\sffamily\bfseries\footnotesize, fill=white, inner sep=1pt, rotate=90]
    at ($(okalt)!0.5!(oküst)$) {taşan yazma};

  % Sağdaki açıklama metinleri
  \node[align=left, text width=48mm, right=14mm of kan.north east, anchor=north west] (not1)
    {Dönüş adresine ulaşmak için önce KANARYA ezilmek zorunda.};
  \node[align=left, text width=48mm, text=iyi, font=\sffamily\bfseries, below=6mm of not1.south west, anchor=north west] (not2)
    {Dönüşte kanarya bozuksa program DURUR.};

  % Dip şeridi
  \node[dip, text width=150mm, below=10mm of don.south -| buf, anchor=north] {%
    Durdurmaz: hedefli yazma · öbek (heap) taşması · bilgi sızıntısı.};
\end{tikzpicture}
\end{document}
